% FILENAME: paper/section-AxialVectorCondensate.tex

\section{Phenomenological Implications: Axial-Vector Condensate}
\label{sec:AxialVectorCondensate}

The majority of this whitepaper is focused on showing that \spinFourC reproduces various aspects of standard physics, but in this section we will explore a consequence of CGD that is drastically different from the Standard Model. For this investigation, we first consider the empty vacuum where the anti-self-dual (ASD) sector's curvature is zero. If the self-dual (SD) sector were structurally equivalent to the ASD sector, then both sectors would be zero and thus the field would trivially commute with itself, meaning all fields would point in the same direction and the universe would only be capable of creating a 1-dimensional object (Equation \ref{eq:kinematicSingleColorDegeneracy}). This would contradict the CGD Axiom \ref{eq:macroscopicVolumeAxiom} asserting that the universe has volume, so we conclude that the universe must have chirality:
\begin{equation}
    \left( \forall \mu, \nu, \ F^{(R)}_{\mu\nu}(x) = 0 \right) \implies \mathcal{A}^{(L)} \neq \mathcal{A}^{(R)}.
    \quad \leanref{CGD.Phenomenology.macroscopicVolumeImpliesChirality}
    \label{eq:macroscopicVolumeImpliesChirality}
\end{equation}
\noindent Furthermore, by looking at the definition of an axial field, $\mathcal{A}_{\text{axial}} = \frac{1}{2}(\mathcal{A}^{(L)} - \mathcal{A}^{(R)})$, we see that since chirality exists throughout the vacuum, a non-zero axial background condensate must exist:
\begin{equation}
    \left( \forall \mu, \nu, \ F^{(R)}_{\mu\nu}(x) = 0 \right) \implies \exists y, \mu \text{ s.t. } \mathcal{A}^{\text{axial}}_\mu(y) \neq 0.
    \quad \leanref{CGD.Phenomenology.macroscopicVolumeImpliesAxialCondensate}
    \label{eq:macroscopicVolumeImpliesAxialCondensate}
\end{equation}

The SD and ASD chiral sectors decompose into symmetric vector $(\mathcal{V})$ and axial $(\mathcal{A}^{\text{axial}})$ background fields:
\begin{align}
    \mathcal{A}^{(L)}_\mu(x) &= \mathcal{V}_\mu(x) + \mathcal{A}^{\text{axial}}_\mu(x)
    \quad \leanref{CGD.Phenomenology.kinematicLeftHandedCoupling}
    \label{eq:kinematicLeftHandedCoupling} \\
    \mathcal{A}^{(R)}_\mu(x) &= \mathcal{V}_\mu(x) - \mathcal{A}^{\text{axial}}_\mu(x) .
    \quad \leanref{CGD.Phenomenology.kinematicRightHandedCoupling}
    \label{eq:kinematicRightHandedCoupling}
\end{align}
\noindent This property is the source of the parity asymmetry of the vacuum.

When we evaluate the gauge-covariant derivative of the fermion state $\Psi$ in the SD sector, we find that the particle interacts with both the vector and axial background fields:
\begin{equation}
    D^{(L)}_\mu \Psi = \partial_\mu \Psi + \mathcal{V}_\mu \Psi + \mathcal{A}^{\text{axial}}_\mu \Psi .
    \quad \leanref{CGD.Phenomenology.kinematicLeftHandedPhaseShift}
    \label{eq:kinematicLeftHandedPhaseShift}
\end{equation}
\noindent This shows why the universe cares about the chirality of particles. In contrast, the Standard Model added the bespoke term $\frac{1}{2}(1 - \gamma^5)$ \citep{feynman1958theory} into the weak interaction current portion of its Lagrangian to describe the chirality discovered in \citet{wu1957experimental}.

We can further characterize the condensate by observing that both the SD and ASD Lie algebras are restricted to \slTwoCAlg, so they are traceless matrices. Since a traceless matrix corresponds with a vector field, the condensate is isovector
\begin{equation}
    \mathrm{Tr}(\mathcal{A}^{\mathrm{axial}}_\mu(x)) = 0 .
    \quad \leanref{CGD.Phenomenology.axialIsIsovector}
    \label{eq:axialIsIsovector}
\end{equation}
\noindent We can also consider what happens when we invert the parity of the field $\mathcal{P}(\mathcal{A}_{\text{axial}}) = \mathcal{P}\left[ \frac{1}{2}(\mathcal{A}^{(L)} - \mathcal{A}^{(R)}) \right]$. This simplifies to $\mathcal{P}(\mathcal{A}_{\text{axial}}) = - \mathcal{A}_{\text{axial}}$, which means
\begin{equation}
    \mathcal{A}^{\mathrm{axial}}_\mu(\mathcal{P}x) = - \mathcal{A}^{\mathrm{axial}}_\mu(x).
    \quad \leanref{CGD.Phenomenology.axialIsParityOdd}
    \label{eq:axialIsParityOdd}
\end{equation}
\noindent All together, we can characterize the quantum numbers of the condensate as $I^G(J^{PC}) = 1^{-}(1^{++})$.

This mathematical necessity of CGD is the crux of why CGD makes novel predictions; in the Standard Model, no axial-vector condensate exists.

At this point, we will stray away from theoretical formalization to look at phenomenological implication of an isovector axial-vector condensate. Only one particle family matches the axial-vector isovector quantum numbers: the $a_1$ meson \citep{pdg2024}, so if we are going to guess that the characteristic particle of the CGD condensate has already been discovered, this is our only choice.

Now that we have characterized the CGD condensate, let us consider how an $a_1$ condensate would impact the proton radius puzzle.

\citet{pohl2010size} showed that measuring the proton radius using muonic hydrogen differed by 5.0 standard deviations compared to the atomic hydrogen measurements. \citet{miller2026proton} argues that newer atomic hydrogen experiments have resolved the discrepancy because \citet{maisenbacher2026sub} used the 2S-6P transition of atomic hydrogen to obtain the measurement of $r_p = 0.8406(15)$ fm, which is consistent with muonic hydrogen measurements. However, this brings up the question of why older atomic hydrogen experiments measured a different value. Physicists have not identified any technical problems for all of the older atomic hydrogen experiments, so in order to declare the proton radius solved, the Standard Model must ignore this data.

Based on QFT, if we bombard a proton with an electron, we need to compute the probability amplitude that the two particles will interact. In the presence of an axial-vector condensate, this amplitude is the product of the bare interaction $G_{Dipole}(Q^2, r_{core})$ (the probability of interacting due to the electron interacting with the proton core) and the vacuum corrections $1 - F_{Axial}(Q^2, \kappa_E)$ (which is the influence of the $a_1$ condensate):
\begin{equation}
    G_E(Q^2) = G_{Dipole}(Q^2, r_{core}) \times ( 1 - F_{Axial}(Q^2, \kappa_E) ).
\end{equation}
\noindent We can then use the standard phenomenological model of Vector Meson Dominance (VMD) \citep{bodek2008vector}, where we observe that when a photon interacts with a proton, it often acts like a short-lived meson. We describe the core interaction using
\begin{equation}
    G_{Dipole}(Q^2, r) = \left( 1 + \frac{Q^2}{12 / (r / \hbar c)^2} \right)^{-2}.
\end{equation}

This leaves us needing to describe the influence of the axial-vector condensate in the $F_{Axial}(Q^2, \kappa_E)$ equation. This is the product of three components: the massive propagator $Q^2 / (Q^2 + m_{a_1}^2)$ associated with the probability amplitude that a virtual photon will resolve into an $a_1$ meson state \citep{guichon1983axial}, the loop screening term $\ln(1 + Q^2 / (m_{a_1}^2 - m_\rho^2))$ which describes the mass gap between the vector environment of the proton and the axial-vector state of the $a_1$ condensate \citep{coleman1973radiative, gross1973asymptotically}, and finally the Weinberg mixing angle \citep{weinberg1967precise} which describes the strength at which the vector and axial-vector environments mix. Putting this together we have
\begin{equation}
    F_{Axial}(Q^2, \kappa_E) = \frac{Q^2}{Q^2 + m_{a_1}^2} \times \ln\left( 1 + \frac{Q^2}{m_{a_1}^2 - m_\rho^2} \right) \times \left( \frac{m_\rho}{m_{a_1}} \right)^2 .
\end{equation}

Notice that the only free parameter in this model is $r$. However, this hides uncertainty around the $a_1$ meson particle family, because \citet{pdg2024} lists multiple resonances for the $a_1$ and $\rho$ mesons. In order to address this issue, in Figure \ref{fig:proton_radius} and Table \ref{tab:proton_radius_fits} we analyze an additional model where we treat $m_{a_1}$ and $m_\rho$ as free parameters.

Another nuance of this proton form factor model is that it requires a relatively high momentum transfer $Q^2$ in order for $F_{Axial}(Q^2, \kappa_E)$ to have any value other than 0. In order to test this model in a regime where the effects of the axial-vector condensate are visible, we will need to find a dataset with a high momentum transfer. For this, we will use the Jefferson Lab high GeV$^2$ polarization transfer experiments described in \citet{punjabi2005proton, gayou2001measurements, puckett2012final, puckett2010recoil} and aggregated by \citet{ye2018proton}.

\begin{figure}[htbp]
    \centering
    \includegraphics[width=\linewidth]{proton-radius/jlab-hall-C-polarization-transfer.pdf}
    \caption{The proton radius factor form factor models. \citet{ye2018proton} is a global parameterization that represents a best-fit for atomic hydrogen based experiments. The CGD 1-Param model leaves only the proton radius $r_p$ as an unknown value, while the CGD 3-Param model additionally treats the meson masses $m_{a_1}$ and $m_\rho$ as unknown values. The statistical analysis of these models is shown in Table \ref{tab:proton_radius_fits}.}
    \label{fig:proton_radius}
\end{figure}

\begin{table}[h]
\centering
\begin{tabular}{lcccc}
\toprule
\textbf{Model} & \textbf{Parameters} & \textbf{Extracted $R_p$ [fm]} & \textbf{Total $\chi^2$} \\
\midrule
CGD Topology (Fixed Mesons) & 1 & 0.8373 & 20.7 \\
CGD Topology (Free Mesons) & 3 & 0.8440 & 17.8 \\
Ye 2018 Global & 13 & 0.879* & 36.0 \\
\bottomrule
\end{tabular}
\caption{The statistical analysis for the models shown in Figure \ref{fig:proton_radius}. The CGD Topology (Free Mesons) model extracted the values $m_{a_1} = 1.2002$ GeV, $m_\rho = 0.7497$ GeV, and $r_p = 0.8440$ fm. The CGD Topology (Fixed Mesons) used the values $m_{a_1} = 1.209$ GeV, $m_\rho = 0.775$ reported in \citet{pdg2024}, and extracted the value $r_p = 0.8373$ fm.}
\label{tab:proton_radius_fits}
\end{table}

It's important to note that despite the statistical comparisons shown in Table \ref{tab:proton_radius_fits}, we can only really reach one conclusion from this: CGD's axial condensate requirement is not easily falsified. We only analyzed a small subset of the available proton radius data, and we deliberately avoided looking at low GeV$^2$ momentum transfers. Furthermore, while our comparison against the Ye 2018 Global model is the best comparison we could find as a proxy for the Standard Model, it is somewhat disingenuous. \citet{ye2018proton} built this parameterization from all of the data available at the time, not just the high GeV$^2$ polarization transfer from JLab Hall C.

